\documentclass[a4paper,11pt]{article}
\usepackage{lmodern}
\usepackage[T1]{fontenc}
\usepackage[utf8]{inputenc}
\usepackage{amsmath,amssymb,amsthm}
\usepackage{hyperref}
\usepackage[margin=25mm]{geometry}

\newtheorem{theorem}{Theorem}[section]
\newtheorem{proposition}[theorem]{Proposition}
\newtheorem{lemma}[theorem]{Lemma}
\newtheorem{corollary}[theorem]{Corollary}
\newtheorem{definition}[theorem]{Definition}
\newtheorem{remark}[theorem]{Remark}
\newtheorem{observation}[theorem]{Observation}

\providecommand{\tightlist}{\setlength{\itemsep}{0pt}\setlength{\parskip}{0pt}}

\begin{document}

\section{Classification and Structural Analysis of the 19 Arbitrary
Parameters of the Standard
Model}\label{classification-and-structural-analysis-of-the-19-arbitrary-parameters-of-the-standard-model}

\subsection{--- Evaluation Based on Structural Requirements for the
Theory of Everything
---}\label{evaluation-based-on-structural-requirements-for-the-theory-of-everything}

\textbf{Author:} Noriaki Kihara (WF System Co., Ltd.)

\textbf{Date:} April 2026

\textbf{DOI:}
\href{https://doi.org/10.5281/zenodo.19604965}{10.5281/zenodo.19604965}

\textbf{Prerequisite Paper:}
\href{https://doi.org/10.5281/zenodo.19601592}{Formulation of Structural
Requirements That a Theory of Everything Must Satisfy}

\begin{center}\rule{0.5\linewidth}{0.5pt}\end{center}

\subsection{0. Purpose of This Paper}\label{purpose-of-this-paper}

This paper classifies and analyzes the 19 arbitrary parameters (free
parameters) contained in the Standard Model of particle physics from the
perspective of structural requirements formulated in a prior work.

Specifically, we address the following questions:

\begin{enumerate}
\def\labelenumi{\arabic{enumi}.}
\tightlist
\item
  Is each parameter structurally derivable, or is it truly arbitrary?
\item
  Which parameters should a Theory of Everything (TOE) internally
  generate (= TOE KPIs)?
\item
  What are the dependency relationships among the parameters?
\end{enumerate}

\begin{center}\rule{0.5\linewidth}{0.5pt}\end{center}

\subsection{1. Parameter Catalog}\label{parameter-catalog}

We present the full picture of parameters involved in the fundamental
theories of modern physics (Standard Model + General Relativity). By
convention, the Standard Model counts ``19 arbitrary parameters,'' but
for evaluating a Theory of Everything, we include all parameters as
verification targets, including those treated as ``fundamental
constants.''

The values of each parameter are those in the low-energy regime
currently confirmed by experiment (roughly the eV to TeV scale, i.e.,
the range accessible to atomic and particle experiments).

\subsubsection{1.0 Fundamental Physical Constants (Currently Treated as
``Constants'')}\label{fundamental-physical-constants-currently-treated-as-constants}

The following are treated as constants in current physics and are
eliminated in natural units (\(c = \hbar = k_B = 1\)), so they are not
counted among the ``19 arbitrary parameters.'' However, among candidates
for grand unified theories and quantum gravity, there exist hypotheses
that these vary over cosmic history or are derivable from more
fundamental quantities. We include them as verification targets for a
Theory of Everything.

These constants are classified by the dimensional composition of their
constituent quantities. This classification clarifies which layer of
physics each constant belongs to.

\paragraph{\texorpdfstring{Group I: Pure Spacetime
(\(L, T\))}{Group I: Pure Spacetime (L, T)}}\label{group-i-pure-spacetime-l-t}

{\def\LTcaptype{none} % do not increment counter
\begin{longtable}[]{@{}
  >{\raggedright\arraybackslash}p{(\linewidth - 10\tabcolsep) * \real{0.0566}}
  >{\raggedright\arraybackslash}p{(\linewidth - 10\tabcolsep) * \real{0.2075}}
  >{\raggedright\arraybackslash}p{(\linewidth - 10\tabcolsep) * \real{0.1509}}
  >{\raggedright\arraybackslash}p{(\linewidth - 10\tabcolsep) * \real{0.1321}}
  >{\raggedright\arraybackslash}p{(\linewidth - 10\tabcolsep) * \real{0.2264}}
  >{\raggedright\arraybackslash}p{(\linewidth - 10\tabcolsep) * \real{0.2264}}@{}}
\toprule\noalign{}
\begin{minipage}[b]{\linewidth}\raggedright
\#
\end{minipage} & \begin{minipage}[b]{\linewidth}\raggedright
Parameter
\end{minipage} & \begin{minipage}[b]{\linewidth}\raggedright
Symbol
\end{minipage} & \begin{minipage}[b]{\linewidth}\raggedright
Value
\end{minipage} & \begin{minipage}[b]{\linewidth}\raggedright
Dimensions
\end{minipage} & \begin{minipage}[b]{\linewidth}\raggedright
Conversion
\end{minipage} \\
\midrule\noalign{}
\endhead
\bottomrule\noalign{}
\endlastfoot
P1 & Speed of light & c & 299,792,458 m/s (definition) & \([L T^{-1}]\)
& Space \(L\) \(\leftrightarrow\) Time \(T\): \(L = cT\) \\
\end{longtable}
}

\begin{itemize}
\tightlist
\item
  Neither mass, charge, nor temperature is involved. A constant of
  \textbf{pure geometry}.
\item
  The foundation of special relativity and the sole constant that
  unifies spacetime.
\item
  TOE question: Can it be derived from the geometrical structure of
  spacetime? Variable speed of light (VSL) theories hypothesize it may
  vary.
\end{itemize}

\paragraph{\texorpdfstring{Group II: Spacetime + Mass
(\(M, L, T\))}{Group II: Spacetime + Mass (M, L, T)}}\label{group-ii-spacetime-mass-m-l-t}

{\def\LTcaptype{none} % do not increment counter
\begin{longtable}[]{@{}
  >{\raggedright\arraybackslash}p{(\linewidth - 10\tabcolsep) * \real{0.0566}}
  >{\raggedright\arraybackslash}p{(\linewidth - 10\tabcolsep) * \real{0.2075}}
  >{\raggedright\arraybackslash}p{(\linewidth - 10\tabcolsep) * \real{0.1509}}
  >{\raggedright\arraybackslash}p{(\linewidth - 10\tabcolsep) * \real{0.1321}}
  >{\raggedright\arraybackslash}p{(\linewidth - 10\tabcolsep) * \real{0.2264}}
  >{\raggedright\arraybackslash}p{(\linewidth - 10\tabcolsep) * \real{0.2264}}@{}}
\toprule\noalign{}
\begin{minipage}[b]{\linewidth}\raggedright
\#
\end{minipage} & \begin{minipage}[b]{\linewidth}\raggedright
Parameter
\end{minipage} & \begin{minipage}[b]{\linewidth}\raggedright
Symbol
\end{minipage} & \begin{minipage}[b]{\linewidth}\raggedright
Value
\end{minipage} & \begin{minipage}[b]{\linewidth}\raggedright
Dimensions
\end{minipage} & \begin{minipage}[b]{\linewidth}\raggedright
Conversion
\end{minipage} \\
\midrule\noalign{}
\endhead
\bottomrule\noalign{}
\endlastfoot
P2 & Planck constant & \(\hbar\) & 1.055 \(\times\) 10\(^{-34}\)
J\(\cdot\)s (definition) & \([M L^2 T^{-1}]\) (\(M^{+1}\)) & Energy
\(E\) \(\leftrightarrow\) Frequency \(\nu\): \(E = \hbar\omega\) \\
P3 & Gravitational constant & G & 6.674 \(\times\) 10\(^{-11}\)
m\(^3\)/(kg\(\cdot\)s\(^2\)) & \([L^3 M^{-1} T^{-2}]\) (\(M^{-1}\)) &
Mass \(M\) \(\leftrightarrow\) Spacetime curvature \(L^{-2}\):
\(R \sim GM/(c^2 r)\) \\
\end{longtable}
}

\begin{itemize}
\tightlist
\item
  Both are composed of the same three dimensions \(M, L, T\), but
  \textbf{the sign of \(M\) is opposite}.

  \begin{itemize}
  \tightlist
  \item
    \(\hbar\): \(M^{+1}\) --- as mass increases, energy (frequency)
    increases
  \item
    \(G\): \(M^{-1}\) --- as mass increases, the gravitational
    constant's contribution inversely decreases (the curvature radius
    grows)
  \end{itemize}
\item
  This \textbf{opposite behavior with respect to mass} already suggests
  at the level of dimensional analysis the fundamental reason why
  quantum mechanics and general relativity are difficult to unify.
\item
  \(\hbar\) is the foundation of quantum mechanics (discretization),
  \(G\) is the foundation of general relativity (geometrization).
\item
  \(G\) has a \textbf{duality}: it is both a dimensional conversion
  constant and the gravitational coupling constant (strength of the
  force).
\end{itemize}

\paragraph{\texorpdfstring{Group III: Spacetime + Mass + Temperature
(\(M, L, T, \Theta\))}{Group III: Spacetime + Mass + Temperature (M, L, T, \textbackslash Theta)}}\label{group-iii-spacetime-mass-temperature-m-l-t-theta}

{\def\LTcaptype{none} % do not increment counter
\begin{longtable}[]{@{}
  >{\raggedright\arraybackslash}p{(\linewidth - 10\tabcolsep) * \real{0.0566}}
  >{\raggedright\arraybackslash}p{(\linewidth - 10\tabcolsep) * \real{0.2075}}
  >{\raggedright\arraybackslash}p{(\linewidth - 10\tabcolsep) * \real{0.1509}}
  >{\raggedright\arraybackslash}p{(\linewidth - 10\tabcolsep) * \real{0.1321}}
  >{\raggedright\arraybackslash}p{(\linewidth - 10\tabcolsep) * \real{0.2264}}
  >{\raggedright\arraybackslash}p{(\linewidth - 10\tabcolsep) * \real{0.2264}}@{}}
\toprule\noalign{}
\begin{minipage}[b]{\linewidth}\raggedright
\#
\end{minipage} & \begin{minipage}[b]{\linewidth}\raggedright
Parameter
\end{minipage} & \begin{minipage}[b]{\linewidth}\raggedright
Symbol
\end{minipage} & \begin{minipage}[b]{\linewidth}\raggedright
Value
\end{minipage} & \begin{minipage}[b]{\linewidth}\raggedright
Dimensions
\end{minipage} & \begin{minipage}[b]{\linewidth}\raggedright
Conversion
\end{minipage} \\
\midrule\noalign{}
\endhead
\bottomrule\noalign{}
\endlastfoot
P4 & Boltzmann constant & \(k_B\) & 1.381 \(\times\) 10\(^{-23}\) J/K
(definition) & \([M L^2 T^{-2} \Theta^{-1}]\) & Temperature \(\Theta\)
\(\leftrightarrow\) Energy \(E\): \(E = k_B T\) \\
\end{longtable}
}

\begin{itemize}
\tightlist
\item
  Adds the temperature dimension \(\Theta\) to Group II (\(M, L, T\)).
\item
  The constant that bridges statistical mechanics (many-body,
  macroscopic) and microscopic physics.
\item
  Bridges informational entropy and thermodynamic entropy. In an
  information-theoretic TOE, \(k_B = 1\) is natural, and it may not be
  an independent parameter.
\end{itemize}

\paragraph{\texorpdfstring{Group IV: Charge Only
(\(Q\))}{Group IV: Charge Only (Q)}}\label{group-iv-charge-only-q}

{\def\LTcaptype{none} % do not increment counter
\begin{longtable}[]{@{}
  >{\raggedright\arraybackslash}p{(\linewidth - 10\tabcolsep) * \real{0.0566}}
  >{\raggedright\arraybackslash}p{(\linewidth - 10\tabcolsep) * \real{0.2075}}
  >{\raggedright\arraybackslash}p{(\linewidth - 10\tabcolsep) * \real{0.1509}}
  >{\raggedright\arraybackslash}p{(\linewidth - 10\tabcolsep) * \real{0.1321}}
  >{\raggedright\arraybackslash}p{(\linewidth - 10\tabcolsep) * \real{0.2264}}
  >{\raggedright\arraybackslash}p{(\linewidth - 10\tabcolsep) * \real{0.2264}}@{}}
\toprule\noalign{}
\begin{minipage}[b]{\linewidth}\raggedright
\#
\end{minipage} & \begin{minipage}[b]{\linewidth}\raggedright
Parameter
\end{minipage} & \begin{minipage}[b]{\linewidth}\raggedright
Symbol
\end{minipage} & \begin{minipage}[b]{\linewidth}\raggedright
Value
\end{minipage} & \begin{minipage}[b]{\linewidth}\raggedright
Dimensions
\end{minipage} & \begin{minipage}[b]{\linewidth}\raggedright
Conversion
\end{minipage} \\
\midrule\noalign{}
\endhead
\bottomrule\noalign{}
\endlastfoot
P5 & Elementary charge & e & 1.602 \(\times\) 10\(^{-19}\) C
(definition) & \([Q]\) & Minimum unit of charge (standard of
quantization) \\
\end{longtable}
}

\begin{itemize}
\tightlist
\item
  Depends on none of \(L, T, M, \Theta\). \textbf{Intrinsically
  independent of spacetime geometry.}
\item
  Constitutes the electromagnetic coupling constant
  \(\alpha = e^2/(4\pi\epsilon_0\hbar c)\).
\item
  TOE question: Why is charge quantized? By Dirac's quantization
  condition, if magnetic monopoles exist, charge quantization becomes
  inevitable.
\end{itemize}

\paragraph{\texorpdfstring{Group V: Charge + Spacetime + Mass
(\(M, L, T, Q\))}{Group V: Charge + Spacetime + Mass (M, L, T, Q)}}\label{group-v-charge-spacetime-mass-m-l-t-q}

{\def\LTcaptype{none} % do not increment counter
\begin{longtable}[]{@{}
  >{\raggedright\arraybackslash}p{(\linewidth - 10\tabcolsep) * \real{0.0566}}
  >{\raggedright\arraybackslash}p{(\linewidth - 10\tabcolsep) * \real{0.2075}}
  >{\raggedright\arraybackslash}p{(\linewidth - 10\tabcolsep) * \real{0.1509}}
  >{\raggedright\arraybackslash}p{(\linewidth - 10\tabcolsep) * \real{0.1321}}
  >{\raggedright\arraybackslash}p{(\linewidth - 10\tabcolsep) * \real{0.2264}}
  >{\raggedright\arraybackslash}p{(\linewidth - 10\tabcolsep) * \real{0.2264}}@{}}
\toprule\noalign{}
\begin{minipage}[b]{\linewidth}\raggedright
\#
\end{minipage} & \begin{minipage}[b]{\linewidth}\raggedright
Parameter
\end{minipage} & \begin{minipage}[b]{\linewidth}\raggedright
Symbol
\end{minipage} & \begin{minipage}[b]{\linewidth}\raggedright
Value
\end{minipage} & \begin{minipage}[b]{\linewidth}\raggedright
Dimensions
\end{minipage} & \begin{minipage}[b]{\linewidth}\raggedright
Conversion
\end{minipage} \\
\midrule\noalign{}
\endhead
\bottomrule\noalign{}
\endlastfoot
P6 & Vacuum permeability/permittivity & \(\mu_0\), \(\epsilon_0\) &
\(c = 1/\sqrt{\mu_0 \epsilon_0}\) & \(\mu_0\): \([M L Q^{-2}]\),
\(\epsilon_0\): \([Q^2 T^2 M^{-1} L^{-3}]\) & Connects charge \(Q\) to
spacetime \((L,T)\) and mass \(M\) \\
\end{longtable}
}

\begin{itemize}
\tightlist
\item
  Connects the charge \(Q\) of Group IV to the spacetime and mass of
  Groups I--II.
\item
  Dependent on \(c\) and thus not independent, but plays the structural
  role of \textbf{determining how isolated charge connects to
  spacetime}.
\item
  If the vacuum is not ``empty'' (has quantum structure), the meaning of
  these constants may change.
\end{itemize}

\paragraph{Dimensional Hierarchy}\label{dimensional-hierarchy}

From the above classification, the fundamental dimensions have the
following hierarchy:

\[\underbrace{L, T}_{\text{Spacetime (geometry)}} \xrightarrow{+M} \underbrace{M, L, T}_{\text{Matter (QM / gravity)}} \xrightarrow{+\Theta} \underbrace{M, L, T, \Theta}_{\text{Statistics (information)}} \]

\[\underbrace{Q}_{\text{Charge (isolated)}} \xrightarrow{\mu_0, \epsilon_0} \underbrace{M, L, T, Q}_{\text{Electromagnetism}}\]

In natural units (\(c = \hbar = k_B = G = 1\)), the distinctions among
\(L, T, M, \Theta\) are eliminated, and all physical quantities reduce
to dimensionless pure numbers. The charge dimension \(Q\) is separately
eliminated by setting \(e = 1\).

\textbf{The ``numbers'' that remain are the true arbitrary parameters
that a Theory of Everything must explain.}

\textbf{Note:} The 2019 SI unit reform made c, \(\hbar\), \(k_B\), e all
\textbf{exact values} (definitions). This does not mean ``their values
were discovered'' but rather ``the units were aligned to these
constants.'' The physical origin of these constants remains unresolved.

\subsubsection{1.1 Independence Analysis of the ``19 Arbitrary
Parameters''}\label{independence-analysis-of-the-19-arbitrary-parameters}

Below we enumerate the 19 parameters conventionally called ``arbitrary''
in the Standard Model, plus gravity, cosmology, and neutrino parameters,
and analyze whether each is \textbf{a truly independent parameter or a
dependent quantity derivable from the fundamental constants of §1.0 or
other parameters}.

\paragraph{A. Coupling Constants of the Forces (Conventionally
4)}\label{a.-coupling-constants-of-the-forces-conventionally-4}

\textbf{\#1 Electromagnetic coupling constant
\(\alpha \approx 1/137.036\)}

\[\alpha = \frac{e^2}{4\pi\epsilon_0\hbar c}\]

\begin{itemize}
\tightlist
\item
  \textbf{Independent? \(\times\) Dependent.} Merely a combination of
  the fundamental constants P1(c), P2(\(\hbar\)), P4(e),
  P6(\(\epsilon_0\)) from §1.0.
\item
  \(\alpha \approx 1/137\) is called a ``mysterious number,'' but
  substituting the values of e, \(\epsilon_0\), \(\hbar\), c
  automatically yields this value. There is nothing to derive.
\item
  The true question is ``why does the elementary charge \(e\) have this
  value'' --- \(\alpha\) is not the substance of the question.
\item
  \textbf{Should not be counted as an arbitrary parameter.
  Double-counted with \(e\) (P4).}
\end{itemize}

\textbf{\#2 Strong force coupling constant \(\alpha_s \approx 0.12\) (at
91 GeV)}

\[\alpha_s(Q^2) \approx \frac{12\pi}{(33-2N_f)\ln(Q^2/\Lambda_{QCD}^2)} \quad \text{(1-loop approximation)}\]

\begin{itemize}
\tightlist
\item
  \textbf{Independent? \(\triangle\) Conditional.} The value of
  \(\alpha_s\) itself depends on the energy scale \(Q\) and is not a
  constant. The sole independent parameter is
  \(\Lambda_{QCD} \approx 200\) MeV (QCD scale).
\item
  However, this derivation formula is a 1-loop perturbative expansion
  approximation that breaks down at low energies (large distances). No
  exact solution has been found (Millennium Prize Problem).
\item
  The spatial dependence is logarithmic, not \(1/r^2\) (power law).
  Qualitatively different from gravity and electromagnetism.
\item
  \textbf{\(\Lambda_{QCD}\) is recognized as an independent parameter,
  but \(\alpha_s\) itself is a dependent quantity.}
\end{itemize}

\textbf{\#3 Weak force coupling constant \(\alpha_w \approx 1/30\)}

\[\alpha_w = \frac{g^2}{4\pi} , \quad G_F = \frac{\sqrt{2}}{8} \frac{g^2}{M_W^2} = \frac{1}{\sqrt{2}v^2}\]

\begin{itemize}
\tightlist
\item
  \textbf{Independent? \(\triangle\) Conditional.} The Fermi constant
  \(G_F\) is derived from the Higgs vacuum expectation value \(v\)
  (\(G_F = 1/(\sqrt{2}v^2)\)).
\item
  In the electroweak unification theory, the weak and electromagnetic
  forces share the same origin (SU(2)\(\times\)U(1) symmetry) and are
  mixed by the Weinberg angle \(\theta_W\).
\item
  \textbf{\(v\) (Higgs VEV) and \(\theta_W\) are the independent
  parameters. \(\alpha_w\) itself is a dependent quantity.}
\end{itemize}

\textbf{\#4 Gravitational coupling constant
\(\alpha_G \approx 5.9 \times 10^{-39}\)}

\[\alpha_G = \frac{Gm_p^2}{\hbar c}\]

\begin{itemize}
\tightlist
\item
  \textbf{Independent? \(\times\) Dependent.} A combination of P1(c),
  P2(\(\hbar\)), P5(G) from §1.0 and the particle mass \(m_p\).
\item
  Moreover, ``which mass \(m\) to use'' is arbitrary (proton? electron?
  Planck mass?). The definition itself is not unique.
\item
  \textbf{Should not be counted as an arbitrary parameter.
  Double-counted with G (P5) and particle masses.}
\end{itemize}

\textbf{Summary of coupling constants:} Among the ``four coupling
constants,'' only \textbf{\(\Lambda_{QCD}\)} and \textbf{the Weinberg
angle \(\theta_W\)} can be counted as independent parameters. \(\alpha\)
and \(\alpha_G\) are combinations of the fundamental constants from §1.0
and particle masses, and are not independent.

\paragraph{B. Cosmological Constant
(1)}\label{b.-cosmological-constant-1}

\textbf{\#5 Cosmological constant
\(\Lambda \approx 1.1 \times 10^{-52}\) m\(^{-2}\)}

\begin{itemize}
\tightlist
\item
  \textbf{Independent? \(\checkmark\) Independent.} No method of
  deriving it from other parameters is known.
\item
  The theoretical prediction from quantum field vacuum zero-point energy
  deviates from the observed value by 120 orders of magnitude.
\item
  Its dimension is \([L^{-2}]\) (curvature). It determines the
  geometrical properties of spacetime.
\item
  \textbf{A truly independent arbitrary parameter.}
\end{itemize}

\paragraph{C. Particle Masses (9)}\label{c.-particle-masses-9}

In the Standard Model, particle masses are determined through Yukawa
coupling constants \(y_i\) with the Higgs field via
\(m_i = y_i \cdot v / \sqrt{2}\).

{\def\LTcaptype{none} % do not increment counter
\begin{longtable}[]{@{}
  >{\raggedright\arraybackslash}p{(\linewidth - 8\tabcolsep) * \real{0.0417}}
  >{\raggedright\arraybackslash}p{(\linewidth - 8\tabcolsep) * \real{0.1528}}
  >{\raggedright\arraybackslash}p{(\linewidth - 8\tabcolsep) * \real{0.0833}}
  >{\raggedright\arraybackslash}p{(\linewidth - 8\tabcolsep) * \real{0.5417}}
  >{\raggedright\arraybackslash}p{(\linewidth - 8\tabcolsep) * \real{0.1806}}@{}}
\toprule\noalign{}
\begin{minipage}[b]{\linewidth}\raggedright
\#
\end{minipage} & \begin{minipage}[b]{\linewidth}\raggedright
Parameter
\end{minipage} & \begin{minipage}[b]{\linewidth}\raggedright
Mass
\end{minipage} & \begin{minipage}[b]{\linewidth}\raggedright
Yukawa coupling \(y_i = m_i\sqrt{2}/v\)
\end{minipage} & \begin{minipage}[b]{\linewidth}\raggedright
Independent?
\end{minipage} \\
\midrule\noalign{}
\endhead
\bottomrule\noalign{}
\endlastfoot
6 & Electron \(m_e\) & 0.511 MeV & \(y_e \approx 2.9 \times 10^{-6}\) &
\(\triangle\) \(y_e\) is independent. \(m_e\) is \(y_e \times v\) \\
7 & Muon \(m_\mu\) & 105.7 MeV & \(y_\mu \approx 6.1 \times 10^{-4}\) &
Same \\
8 & Tau \(m_\tau\) & 1776.9 MeV & \(y_\tau \approx 1.0 \times 10^{-2}\)
& Same \\
9 & Up \(m_u\) & \(\approx\) 2.2 MeV &
\(y_u \approx 1.3 \times 10^{-5}\) & Same \\
10 & Down \(m_d\) & \(\approx\) 4.7 MeV &
\(y_d \approx 2.7 \times 10^{-5}\) & Same \\
11 & Charm \(m_c\) & \(\approx\) 1275 MeV &
\(y_c \approx 7.3 \times 10^{-3}\) & Same \\
12 & Strange \(m_s\) & \(\approx\) 95 MeV &
\(y_s \approx 5.5 \times 10^{-4}\) & Same \\
13 & Top \(m_t\) & \(\approx\) 173200 MeV & \(y_t \approx 1.0\) (nearly
1) & Same \\
14 & Bottom \(m_b\) & \(\approx\) 4180 MeV &
\(y_b \approx 2.4 \times 10^{-2}\) & Same \\
\end{longtable}
}

\begin{itemize}
\tightlist
\item
  \textbf{Independent? \(\triangle\) Conditional.} The masses \(m_i\)
  themselves are dependent quantities derived via
  \(y_i \times v/\sqrt{2}\). The truly independent quantities are the
  Yukawa coupling constants \(y_i\) (9) and the Higgs VEV \(v\) (1).
\item
  Notably, \(y_t \approx 1\) (top quark) is the only ``natural'' value,
  while the remaining 8 are unnaturally small at \(10^{-6}\) to
  \(10^{-2}\). \textbf{Why the Yukawa couplings exhibit such hierarchy
  is unsolved.}
\item
  \textbf{The 9 masses decompose into 9 Yukawa coupling constants
  \(y_i\) + \(v\). They should not be double-counted as masses.}
\end{itemize}

\paragraph{D. CKM Matrix Parameters
(4)}\label{d.-ckm-matrix-parameters-4}

{\def\LTcaptype{none} % do not increment counter
\begin{longtable}[]{@{}
  >{\raggedright\arraybackslash}p{(\linewidth - 6\tabcolsep) * \real{0.0882}}
  >{\raggedright\arraybackslash}p{(\linewidth - 6\tabcolsep) * \real{0.3235}}
  >{\raggedright\arraybackslash}p{(\linewidth - 6\tabcolsep) * \real{0.2059}}
  >{\raggedright\arraybackslash}p{(\linewidth - 6\tabcolsep) * \real{0.3824}}@{}}
\toprule\noalign{}
\begin{minipage}[b]{\linewidth}\raggedright
\#
\end{minipage} & \begin{minipage}[b]{\linewidth}\raggedright
Parameter
\end{minipage} & \begin{minipage}[b]{\linewidth}\raggedright
Value
\end{minipage} & \begin{minipage}[b]{\linewidth}\raggedright
Independent?
\end{minipage} \\
\midrule\noalign{}
\endhead
\bottomrule\noalign{}
\endlastfoot
15 & Mixing angle \(\theta_{12}\) & \(\approx\) 13.0\(^\circ\) (Cabibbo
angle) & \(\checkmark\) Independent \\
16 & Mixing angle \(\theta_{23}\) & \(\approx\) 2.4\(^\circ\) &
\(\checkmark\) Independent \\
17 & Mixing angle \(\theta_{13}\) & \(\approx\) 0.2\(^\circ\) &
\(\checkmark\) Independent \\
18 & CP phase \(\delta_{CKM}\) & \(\approx\) 1.2 rad & \(\checkmark\)
Independent \\
\end{longtable}
}

\begin{itemize}
\tightlist
\item
  \textbf{Independent? \(\checkmark\) Independent.} No method of
  derivation from other parameters is known.
\item
  However, if the generation structure of quarks (why 3 generations) is
  elucidated, they might be determined geometrically or
  group-theoretically.
\item
  \textbf{Parameters specific to the weak force. Unnecessary if the weak
  force did not exist.}
\end{itemize}

\paragraph{E. Higgs Field Parameters
(2)}\label{e.-higgs-field-parameters-2}

\textbf{\#19 Higgs boson mass \(m_H \approx 125.1\) GeV}

\[m_H = \sqrt{2\lambda} \cdot v\]

\begin{itemize}
\tightlist
\item
  \textbf{Independent? \(\triangle\) Conditional.} Derived from the
  Higgs self-coupling constant \(\lambda\) and VEV \(v\). The truly
  independent quantities are \(\lambda\) and \(v\).
\item
  Treating \(m_H\) as an ``arbitrary parameter'' merely packages
  \(\lambda \times v\) into a single number.
\end{itemize}

\textbf{\#20 Vacuum expectation value \(v \approx 246.2\) GeV}

\[v = (\sqrt{2}G_F)^{-1/2} = \frac{2M_W}{g}\]

\begin{itemize}
\tightlist
\item
  \textbf{Independent? \(\checkmark\) Independent.} Determines the scale
  of electroweak symmetry breaking. The reference for all particle mass
  scales.
\item
  The range of the weak force (W/Z boson masses) is also determined by
  this value.
\item
  \textbf{One of the most important independent parameters.}
\end{itemize}

\paragraph{F. QCD Vacuum Parameter (1)}\label{f.-qcd-vacuum-parameter-1}

\textbf{\#21 QCD vacuum angle \(\theta_{QCD} < 10^{-10}\)}

\begin{itemize}
\tightlist
\item
  \textbf{Independent? \(\checkmark\) Independent (but the value is
  nearly zero).}
\item
  The experimental upper bound is \(< 10^{-10}\), and it may be
  physically zero.
\item
  If zero, some symmetry (Peccei-Quinn symmetry) underlies it, making it
  not an arbitrary parameter but \textbf{structurally zero}.
\item
  \textbf{Whether it is an ``arbitrary parameter'' or ``structurally
  zero'' is undetermined.}
\end{itemize}

\paragraph{G. Neutrino Parameters (7+)}\label{g.-neutrino-parameters-7}

{\def\LTcaptype{none} % do not increment counter
\begin{longtable}[]{@{}
  >{\raggedright\arraybackslash}p{(\linewidth - 6\tabcolsep) * \real{0.0882}}
  >{\raggedright\arraybackslash}p{(\linewidth - 6\tabcolsep) * \real{0.3235}}
  >{\raggedright\arraybackslash}p{(\linewidth - 6\tabcolsep) * \real{0.2059}}
  >{\raggedright\arraybackslash}p{(\linewidth - 6\tabcolsep) * \real{0.3824}}@{}}
\toprule\noalign{}
\begin{minipage}[b]{\linewidth}\raggedright
\#
\end{minipage} & \begin{minipage}[b]{\linewidth}\raggedright
Parameter
\end{minipage} & \begin{minipage}[b]{\linewidth}\raggedright
Value
\end{minipage} & \begin{minipage}[b]{\linewidth}\raggedright
Independent?
\end{minipage} \\
\midrule\noalign{}
\endhead
\bottomrule\noalign{}
\endlastfoot
22 & \(\Delta m^2_{21}\) & \(\approx 7.5 \times 10^{-5}\) eV\(^2\) &
\(\checkmark\) Independent (squared mass difference; absolute mass
unmeasured) \\
23 & \(\Delta m^2_{32}\) & \(\approx 2.5 \times 10^{-3}\) eV\(^2\) &
\(\checkmark\) Independent \\
24 & PMNS \(\theta_{12}\) & \(\approx\) 33.4\(^\circ\) & \(\checkmark\)
Independent \\
25 & PMNS \(\theta_{23}\) & \(\approx\) 49\(^\circ\) & \(\checkmark\)
Independent \\
26 & PMNS \(\theta_{13}\) & \(\approx\) 8.5\(^\circ\) & \(\checkmark\)
Independent \\
27 & PMNS \(\delta\) & Undetermined & \(\checkmark\) Independent
(unmeasured) \\
28 & Majorana phases \(\alpha_1\), \(\alpha_2\) & Unmeasured & ? Whether
they exist is also undetermined \\
\end{longtable}
}

\subsubsection{1.8 Summary of Independence
Analysis}\label{summary-of-independence-analysis}

Through the above analysis, the reality of the ``19 arbitrary
parameters'' is as follows.

\paragraph{Dependent Quantities (Derivable from Others; Should Not Be
Counted as Independent
Parameters)}\label{dependent-quantities-derivable-from-others-should-not-be-counted-as-independent-parameters}

{\def\LTcaptype{none} % do not increment counter
\begin{longtable}[]{@{}
  >{\raggedright\arraybackslash}p{(\linewidth - 4\tabcolsep) * \real{0.2245}}
  >{\raggedright\arraybackslash}p{(\linewidth - 4\tabcolsep) * \real{0.3878}}
  >{\raggedright\arraybackslash}p{(\linewidth - 4\tabcolsep) * \real{0.3878}}@{}}
\toprule\noalign{}
\begin{minipage}[b]{\linewidth}\raggedright
Parameter
\end{minipage} & \begin{minipage}[b]{\linewidth}\raggedright
Derivation Formula
\end{minipage} & \begin{minipage}[b]{\linewidth}\raggedright
Double-Counted With
\end{minipage} \\
\midrule\noalign{}
\endhead
\bottomrule\noalign{}
\endlastfoot
\(\alpha \approx 1/137\) & \(e^2/(4\pi\epsilon_0\hbar c)\) & P1(c),
P2(\(\hbar\)), P4(e), P6(\(\epsilon_0\)) \\
\(\alpha_G \approx 10^{-39}\) & \(Gm^2/(\hbar c)\) & P1(c),
P2(\(\hbar\)), P5(G) + particle mass. Moreover, \(m\) is arbitrary,
making the definition non-unique \\
\(\alpha_s\) (at specific scale) &
\(12\pi/((33-2N_f)\ln(Q^2/\Lambda_{QCD}^2))\) & Reduces to
\(\Lambda_{QCD}\) \\
\(\alpha_w\) & \(g^2/(4\pi)\) & Reduces to \(v\) and \(\theta_W\) \\
9 particle masses & \(m_i = y_i \cdot v/\sqrt{2}\) & Reduces to Yukawa
couplings \(y_i\) + \(v\) \\
Higgs mass \(m_H\) & \(m_H = \sqrt{2\lambda} \cdot v\) & Reduces to
Higgs self-coupling \(\lambda\) + \(v\) \\
\end{longtable}
}

\paragraph{Truly Independent
Parameters}\label{truly-independent-parameters}

{\def\LTcaptype{none} % do not increment counter
\begin{longtable}[]{@{}
  >{\raggedright\arraybackslash}p{(\linewidth - 4\tabcolsep) * \real{0.3571}}
  >{\raggedright\arraybackslash}p{(\linewidth - 4\tabcolsep) * \real{0.3929}}
  >{\raggedright\arraybackslash}p{(\linewidth - 4\tabcolsep) * \real{0.2500}}@{}}
\toprule\noalign{}
\begin{minipage}[b]{\linewidth}\raggedright
Category
\end{minipage} & \begin{minipage}[b]{\linewidth}\raggedright
Parameters
\end{minipage} & \begin{minipage}[b]{\linewidth}\raggedright
Count
\end{minipage} \\
\midrule\noalign{}
\endhead
\bottomrule\noalign{}
\endlastfoot
Fundamental physical constants & c, \(\hbar\), \(k_B\), e, G,
(\(\mu_0\epsilon_0\)) & 6 (\(\mu_0\epsilon_0\) dependent on c) \\
Cosmological constant & \(\Lambda\) & 1 \\
Strong force scale & \(\Lambda_{QCD}\) & 1 \\
Electroweak mixing & Weinberg angle \(\theta_W\) & 1 \\
Yukawa couplings & \(y_e, y_\mu, y_\tau, y_u, y_d, y_c, y_s, y_t, y_b\)
& 9 \\
Higgs & Self-coupling \(\lambda\), vacuum expectation value \(v\) & 2 \\
CKM matrix & \(\theta_{12}, \theta_{23}, \theta_{13}, \delta_{CKM}\) &
4 \\
QCD vacuum & \(\theta_{QCD}\) & 1 (possibly structurally zero) \\
Neutrino &
\(\Delta m^2_{21}, \Delta m^2_{32}, \theta_{12}, \theta_{23}, \theta_{13}, \delta, (\alpha_1, \alpha_2)\)
& 7+ \\
\textbf{Total} & & \textbf{\$\approx\$32} (5 fundamental constants +
\$\approx\$27 arbitrary parameters) \\
\end{longtable}
}

\paragraph{Conclusion of This Section
(Provisional)}\label{conclusion-of-this-section-provisional}

The conventional count of ``19 arbitrary parameters'' is inaccurate in
the following respects:

\begin{enumerate}
\def\labelenumi{\arabic{enumi}.}
\tightlist
\item
  \textbf{Double-counting:} \(\alpha\), \(\alpha_G\), \(\alpha_w\),
  particle masses, and \(m_H\) are combinations of other parameters and
  should not be counted as independent.
\item
  \textbf{Under-counting:} Gravity (G, \(\Lambda\)), neutrinos (7+), and
  hidden parameters (\(\Lambda_{QCD}\), \(\theta_W\), Yukawa couplings,
  \(\lambda\)) are not included.
\item
  \textbf{Exclusion of fundamental constants:} c, \(\hbar\), G, \(k_B\),
  e are excluded as ``constants,'' but the origin of their values is
  unresolved.
\end{enumerate}

As a KPI for the Theory of Everything, we evaluate the following for
these approximately 32 parameters:

\[\text{KPI} = \frac{\text{Number of parameters whose theoretically derived values match experiment}}{\text{Total number of parameters under verification ($\approx$32)}}\]

The ideal is to reproduce all \$\approx\$32 from minimal input (0--2).
The current Standard Model's KPI is effectively zero (it merely returns
its inputs).

\begin{center}\rule{0.5\linewidth}{0.5pt}\end{center}

\subsection{2. Structural Analysis of the Coupling
Constants}\label{structural-analysis-of-the-coupling-constants}

As shown in §1.1, the ``four coupling constants'' are all dependent
quantities, but the structure of their derivation formulas conceals
important problems.

\subsubsection{\texorpdfstring{2.1 Electromagnetic Coupling Constant
\(\alpha\)}{2.1 Electromagnetic Coupling Constant \textbackslash alpha}}\label{electromagnetic-coupling-constant-alpha}

\[\alpha = \frac{e^2}{4\pi\epsilon_0\hbar c} \approx \frac{1}{137.036}\]

\textbf{Dimensional analysis of constituents:}

{\def\LTcaptype{none} % do not increment counter
\begin{longtable}[]{@{}
  >{\raggedright\arraybackslash}p{(\linewidth - 4\tabcolsep) * \real{0.2143}}
  >{\raggedright\arraybackslash}p{(\linewidth - 4\tabcolsep) * \real{0.2857}}
  >{\raggedright\arraybackslash}p{(\linewidth - 4\tabcolsep) * \real{0.5000}}@{}}
\toprule\noalign{}
\begin{minipage}[b]{\linewidth}\raggedright
Element
\end{minipage} & \begin{minipage}[b]{\linewidth}\raggedright
Dimensions
\end{minipage} & \begin{minipage}[b]{\linewidth}\raggedright
§1.0 Classification
\end{minipage} \\
\midrule\noalign{}
\endhead
\bottomrule\noalign{}
\endlastfoot
\(e^2\) & \([Q^2]\) & Group IV (Charge) \\
\(\epsilon_0\) & \([Q^2 T^2 M^{-1} L^{-3}]\) & Group V (Charge
\(\leftrightarrow\) Spacetime) \\
\(\hbar\) & \([M L^2 T^{-1}]\) & Group II (Spacetime + Mass) \\
\(c\) & \([L T^{-1}]\) & Group I (Pure Spacetime) \\
\end{longtable}
}

Dimensional verification:
\([Q^2] / ([Q^2 T^2 M^{-1} L^{-3}] \cdot [M L^2 T^{-1}] \cdot [L T^{-1}]) = [1]\)
(dimensionless) \(\checkmark\)

\textbf{Structural meaning:} \(\alpha\) represents the ratio of how
charge (Group IV) couples with spacetime + mass (Groups I, II, V). It
uses 4 of the 5 fundamental constants (c, \(\hbar\), e, \(\epsilon_0\))
to produce a single dimensionless number. \textbf{There is no mystery in
the value of \(\alpha\); if mystery exists, it lies in the value of
\(e\).}

\subsubsection{\texorpdfstring{2.2 Strong Force Coupling Constant
\(\alpha_s\)}{2.2 Strong Force Coupling Constant \textbackslash alpha\_s}}\label{strong-force-coupling-constant-alpha_s}

\[\alpha_s(Q^2) \approx \frac{12\pi}{(33-2N_f)\ln(Q^2/\Lambda_{QCD}^2)} \quad \text{(1-loop approximation)}\]

\textbf{Structural problems:}

\begin{enumerate}
\def\labelenumi{\arabic{enumi}.}
\tightlist
\item
  \textbf{Not a constant.} A function of the energy scale \(Q\)
  (\(\approx\) inverse of spatial resolution \(\hbar/r\)). The term
  ``coupling constant'' is misleading.
\item
  \textbf{Approximate formula.} A result of 1-loop perturbative
  expansion, not an exact solution. Higher-order corrections increase
  complexity.
\item
  \textbf{Breaks down at low energy.} As \(Q \to \Lambda_{QCD}\), the
  logarithm diverges and computation becomes impossible. The confinement
  regime relies on lattice QCD (numerical computation). The existence of
  an exact solution is unproven (Millennium Prize Problem).
\item
  \textbf{Spatial dependence is not geometrical.} Compared to the
  \(1/r^2\) of gravity and electromagnetism (inverse of area in 3D
  space), \(\alpha_s\) is logarithmic. This does not arise naturally
  from dimensional analysis but appears only as a result of quantum
  corrections (virtual particle pair creation and annihilation).
\end{enumerate}

\textbf{The only independent parameter is \(\Lambda_{QCD}\).} However,
\(\Lambda_{QCD}\) plays the role of a ``hidden dimensional conversion
constant that connects color charge (a discrete label) to the spatial
dimension \(L\)'' (unclassified in §1.0).

\subsubsection{\texorpdfstring{2.3 Weak Force Coupling Constant
\(\alpha_w\)}{2.3 Weak Force Coupling Constant \textbackslash alpha\_w}}\label{weak-force-coupling-constant-alpha_w}

\[\alpha_w = \frac{g^2}{4\pi}, \quad \sin^2\theta_W = 1 - \frac{M_W^2}{M_Z^2} \approx 0.231\]

\[G_F = \frac{1}{\sqrt{2}v^2} \approx 1.166 \times 10^{-5} \text{ GeV}^{-2}\]

\textbf{Structural features:}

\begin{enumerate}
\def\labelenumi{\arabic{enumi}.}
\tightlist
\item
  \textbf{Already unified with electromagnetism (electroweak
  unification).} In SU(2)\(\times\)U(1) gauge theory, the weak and
  electromagnetic forces are mixed by the Weinberg angle \(\theta_W\).
  At low energies they appear as separate forces, but above \(\approx\)
  100 GeV they are unified.
\item
  \textbf{Finite range.} Due to the masses of the W/Z bosons
  (\(M_W \approx 80.4\) GeV, \(M_Z \approx 91.2\) GeV), the range of the
  force is limited to \(\sim \hbar/(M_W c) \approx 10^{-18}\) m. It
  follows a Yukawa-type potential \(e^{-M_W r}/r\).
\item
  \textbf{Reduces to the Higgs VEV \(v\).} Since
  \(G_F = 1/(\sqrt{2}v^2)\), the scale of the weak force is entirely
  determined by \(v\).
\end{enumerate}

\textbf{The independent parameters are \(v\) and \(\theta_W\).}
\(\alpha_w\) itself is a dependent quantity.

\subsubsection{\texorpdfstring{2.4 Gravitational Coupling Constant
\(\alpha_G\)}{2.4 Gravitational Coupling Constant \textbackslash alpha\_G}}\label{gravitational-coupling-constant-alpha_g}

\[\alpha_G = \frac{Gm^2}{\hbar c}\]

\textbf{Structural problems:}

\begin{enumerate}
\def\labelenumi{\arabic{enumi}.}
\item
  \textbf{The choice of mass \(m\) is arbitrary.} The other three
  coupling constants are determined solely by properties of the force,
  but \(\alpha_G\) changes depending on ``which particle's mass is
  used.''

  {\def\LTcaptype{none} % do not increment counter
  \begin{longtable}[]{@{}ll@{}}
  \toprule\noalign{}
  Mass choice & Value of \(\alpha_G\) \\
  \midrule\noalign{}
  \endhead
  \bottomrule\noalign{}
  \endlastfoot
  Electron mass \(m_e\) & \(\approx 1.8 \times 10^{-45}\) \\
  Proton mass \(m_p\) & \(\approx 5.9 \times 10^{-39}\) \\
  Planck mass \(m_{Pl}\) & \(= 1\) (by definition) \\
  \end{longtable}
  }

  A ``constant'' whose definition is not even unique should not be
  called a constant.
\item
  \textbf{A combination of §1.0 constants.} G (P5), \(\hbar\) (P2), c
  (P1) are all fundamental constants. \(\alpha_G\) contains no new
  information and is merely a double-counting of particle masses and
  fundamental constants.
\item
  \textbf{``Gravity is weak'' is a question of curvature radius.} The
  smallness of \(\alpha_G\) does not mean ``gravity is weak'' but
  reflects that the curvature radius of spacetime is overwhelmingly
  larger than the curvature radius of electromagnetic acceleration.
  Comparing them at the same scale is itself a category error.
\end{enumerate}

\subsubsection{2.5 Fundamental Problem in Comparing the Four Coupling
Constants}\label{fundamental-problem-in-comparing-the-four-coupling-constants}

Conventional physics lines up the four coupling constants as follows and
poses the ``hierarchy problem'':

{\def\LTcaptype{none} % do not increment counter
\begin{longtable}[]{@{}lll@{}}
\toprule\noalign{}
Force & Coupling constant & Relative ratio \\
\midrule\noalign{}
\endhead
\bottomrule\noalign{}
\endlastfoot
Strong & \(\alpha_s \approx 1\) & 1 \\
Electromagnetic & \(\alpha \approx 1/137\) & \(10^{-2}\) \\
Weak & \(\alpha_w \approx 1/30\) & \(10^{-2}\) \\
Gravity & \(\alpha_G \approx 10^{-39}\) & \(10^{-39}\) \\
\end{longtable}
}

However, our analysis reveals that this comparison is invalid for the
following reasons:

\begin{enumerate}
\def\labelenumi{\arabic{enumi}.}
\tightlist
\item
  \textbf{The objects of comparison are not of the same kind.}
  \(\alpha\) is a combination of fundamental constants, \(\alpha_s\) is
  a scale-dependent function, \(\alpha_w\) is a dependent quantity of
  \(v\) and \(\theta_W\), and \(\alpha_G\) is an indeterminate quantity
  depending on mass selection. Placing these in the same table is itself
  meaningless.
\item
  \textbf{Characteristic scales differ.} Each force operates at its own
  spatial scale (\(10^{-18}\) m to cosmic scale), and there is no basis
  for comparing them at a single energy scale.
\item
  \textbf{Spatial dependencies are qualitatively different.} Gravity and
  electromagnetism follow \(1/r^2\) (geometrical), the strong force
  follows a logarithmic function (quantum corrections), and the weak
  force follows exponential decay (Yukawa type). Comparing the
  ``strength'' of quantities with different mathematical structures by a
  single number is physically meaningless.
\item
  \textbf{The ``hierarchy problem'' is a pseudo-problem.} From points
  1--3 above, the question ``why is gravity alone 39 orders of magnitude
  weaker?'' is itself improperly formulated.
\end{enumerate}

\begin{center}\rule{0.5\linewidth}{0.5pt}\end{center}

\subsection{3. Structural Analysis of Particle
Masses}\label{structural-analysis-of-particle-masses}

\subsubsection{3.1 Hierarchy of Yukawa Coupling
Constants}\label{hierarchy-of-yukawa-coupling-constants}

As shown in §1.1 C, particle masses \(m_i\) are products of the Yukawa
coupling constants \(y_i\) and the Higgs VEV \(v\).

\[m_i = \frac{y_i \cdot v}{\sqrt{2}}\]

The truly independent quantities are \(y_i\), whose hierarchy is as
follows:

{\def\LTcaptype{none} % do not increment counter
\begin{longtable}[]{@{}
  >{\raggedright\arraybackslash}p{(\linewidth - 12\tabcolsep) * \real{0.1714}}
  >{\raggedright\arraybackslash}p{(\linewidth - 12\tabcolsep) * \real{0.1143}}
  >{\raggedright\arraybackslash}p{(\linewidth - 12\tabcolsep) * \real{0.0714}}
  >{\raggedright\arraybackslash}p{(\linewidth - 12\tabcolsep) * \real{0.2286}}
  >{\raggedright\arraybackslash}p{(\linewidth - 12\tabcolsep) * \real{0.0714}}
  >{\raggedright\arraybackslash}p{(\linewidth - 12\tabcolsep) * \real{0.2714}}
  >{\raggedright\arraybackslash}p{(\linewidth - 12\tabcolsep) * \real{0.0714}}@{}}
\toprule\noalign{}
\begin{minipage}[b]{\linewidth}\raggedright
Generation
\end{minipage} & \begin{minipage}[b]{\linewidth}\raggedright
Lepton
\end{minipage} & \begin{minipage}[b]{\linewidth}\raggedright
\(y\)
\end{minipage} & \begin{minipage}[b]{\linewidth}\raggedright
Quark (up-type)
\end{minipage} & \begin{minipage}[b]{\linewidth}\raggedright
\(y\)
\end{minipage} & \begin{minipage}[b]{\linewidth}\raggedright
Quark (down-type)
\end{minipage} & \begin{minipage}[b]{\linewidth}\raggedright
\(y\)
\end{minipage} \\
\midrule\noalign{}
\endhead
\bottomrule\noalign{}
\endlastfoot
1st & \(e\) & 2.9 \(\times\) 10\(^{-6}\) & \(u\) & 1.3 \(\times\)
10\(^{-5}\) & \(d\) & 2.7 \(\times\) 10\(^{-5}\) \\
2nd & \(\mu\) & 6.1 \(\times\) 10\(^{-4}\) & \(c\) & 7.3 \(\times\)
10\(^{-3}\) & \(s\) & 5.5 \(\times\) 10\(^{-4}\) \\
3rd & \(\tau\) & 1.0 \(\times\) 10\(^{-2}\) & \(t\) &
\textbf{\(\approx\) 1} & \(b\) & 2.4 \(\times\) 10\(^{-2}\) \\
\end{longtable}
}

\textbf{Notable structures:}

\begin{enumerate}
\def\labelenumi{\arabic{enumi}.}
\tightlist
\item
  \textbf{Only \(y_t \approx 1\) (top quark) is ``natural.''} The
  remaining 8 are unnaturally small at \(10^{-6}\) to \(10^{-2}\). Why
  \(y_t\) alone matches the Higgs VEV scale is unsolved.
\item
  \textbf{Regularity across generations.} Within the same column
  (leptons, up-type quarks, down-type quarks), the values roughly
  increase by 2--3 orders of magnitude per generation. Whether this is
  coincidence or reflects some structure is unknown.
\item
  \textbf{Are all 9 independent?} If the inter-generation ratios have
  regularity, the independent parameters might be reducible to
  ``reference mass + ratio rule.''
\end{enumerate}

\subsubsection{3.2 The Problem of Mass
Origin}\label{the-problem-of-mass-origin}

The Standard Model states that ``the Higgs mechanism gives mass,'' but
more precisely, the Higgs mechanism provides the \textbf{mechanism} of
mass (spontaneous symmetry breaking) without explaining the
\textbf{values} of masses (Yukawa coupling constants \(y_i\)). Why the 9
\(y_i\) have these values is in principle unanswerable within the
framework of the Standard Model.

\begin{center}\rule{0.5\linewidth}{0.5pt}\end{center}

\subsection{4. Structural Analysis of Mixing Angles and CP
Phases}\label{structural-analysis-of-mixing-angles-and-cp-phases}

\subsubsection{4.1 CKM Matrix (Quarks)}\label{ckm-matrix-quarks}

The CKM matrix is the rotation matrix between the weak force gauge
eigenstates and mass eigenstates. A 3-generation unitary matrix is
completely described by 3 rotation angles and 1 phase.

{\def\LTcaptype{none} % do not increment counter
\begin{longtable}[]{@{}
  >{\raggedright\arraybackslash}p{(\linewidth - 4\tabcolsep) * \real{0.3143}}
  >{\raggedright\arraybackslash}p{(\linewidth - 4\tabcolsep) * \real{0.2000}}
  >{\raggedright\arraybackslash}p{(\linewidth - 4\tabcolsep) * \real{0.4857}}@{}}
\toprule\noalign{}
\begin{minipage}[b]{\linewidth}\raggedright
Parameter
\end{minipage} & \begin{minipage}[b]{\linewidth}\raggedright
Value
\end{minipage} & \begin{minipage}[b]{\linewidth}\raggedright
Physical meaning
\end{minipage} \\
\midrule\noalign{}
\endhead
\bottomrule\noalign{}
\endlastfoot
\(\theta_{12} \approx 13.0^\circ\) & Cabibbo angle & Mixing between
1st--2nd generations. Largest \\
\(\theta_{23} \approx 2.4^\circ\) & & Between 2nd--3rd generations.
Small \\
\(\theta_{13} \approx 0.2^\circ\) & & Between 1st--3rd generations.
Extremely small \\
\(\delta_{CKM} \approx 1.2\) rad & CP phase & Matter--antimatter
asymmetry \\
\end{longtable}
}

\textbf{Structural problems:}

\begin{enumerate}
\def\labelenumi{\arabic{enumi}.}
\tightlist
\item
  \textbf{Why 3 generations?} The CKM matrix is 3\$\times\$3 because
  quarks have 3 generations, but why there are 3 generations is
  unresolved. With 2 generations the CP phase would be unnecessary; with
  4+ generations additional parameters would be required.
\item
  \textbf{Hierarchy of mixing angles.} There is a hierarchy
  \(\theta_{12} > \theta_{23} > \theta_{13}\). This may correlate with
  the Yukawa coupling hierarchy (inter-generation mass ratios).
\item
  \textbf{Specific to the weak force.} The CKM matrix would be
  unnecessary without the weak force. If the existence of the weak force
  itself is derived from structural requirements, the mixing angles may
  also be structurally determined.
\end{enumerate}

\subsubsection{4.2 PMNS Matrix (Neutrinos)}\label{pmns-matrix-neutrinos}

The neutrino version of the CKM matrix. Same structure but vastly
different values.

{\def\LTcaptype{none} % do not increment counter
\begin{longtable}[]{@{}
  >{\raggedright\arraybackslash}p{(\linewidth - 4\tabcolsep) * \real{0.2558}}
  >{\raggedright\arraybackslash}p{(\linewidth - 4\tabcolsep) * \real{0.3023}}
  >{\raggedright\arraybackslash}p{(\linewidth - 4\tabcolsep) * \real{0.4419}}@{}}
\toprule\noalign{}
\begin{minipage}[b]{\linewidth}\raggedright
Parameter
\end{minipage} & \begin{minipage}[b]{\linewidth}\raggedright
CKM (Quarks)
\end{minipage} & \begin{minipage}[b]{\linewidth}\raggedright
PMNS (Neutrinos)
\end{minipage} \\
\midrule\noalign{}
\endhead
\bottomrule\noalign{}
\endlastfoot
\(\theta_{12}\) & 13.0\(^\circ\) (small) & 33.4\(^\circ\) (large) \\
\(\theta_{23}\) & 2.4\(^\circ\) (small) & 49\(^\circ\) (nearly maximal
mixing) \\
\(\theta_{13}\) & 0.2\(^\circ\) (tiny) & 8.5\(^\circ\) (small but
finite) \\
\end{longtable}
}

\textbf{Striking contrast between CKM and PMNS:} Quark mixing is small
(diagonal components dominate) while neutrino mixing is large
(off-diagonal components are prominent). Why the values differ so
greatly despite the same mathematical structure (3\$\times\$3 unitary
matrix) is unsolved.

\begin{center}\rule{0.5\linewidth}{0.5pt}\end{center}

\subsection{5. Structural Analysis of the Higgs Field and Cosmological
Constant}\label{structural-analysis-of-the-higgs-field-and-cosmological-constant}

\subsubsection{5.1 Higgs Field}\label{higgs-field}

\textbf{Independent parameters:} Self-coupling constant
\(\lambda \approx 0.13\), vacuum expectation value \(v \approx 246.2\)
GeV

\[V(\phi) = -\mu^2|\phi|^2 + \lambda|\phi|^4, \quad v = \frac{\mu}{\sqrt{\lambda}}, \quad m_H = \sqrt{2\lambda}\cdot v\]

\textbf{Structural role:}

\begin{itemize}
\tightlist
\item
  \(v\) determines the scale of electroweak symmetry breaking and is the
  most important parameter as the reference for all particle masses.
\item
  \(\lambda\) determines the mass of the Higgs boson itself, but there
  is no theoretical basis for the value \(\approx 0.13\).
\item
  The origins of \(v\) and \(\lambda\) cannot be explained within the
  framework of the Standard Model. ``Why symmetry breaks'' is merely
  assumed structurally.
\end{itemize}

\subsubsection{\texorpdfstring{5.2 Cosmological Constant
\(\Lambda\)}{5.2 Cosmological Constant \textbackslash Lambda}}\label{cosmological-constant-lambda}

\[G_{\mu\nu} + \Lambda g_{\mu\nu} = \frac{8\pi G}{c^4} T_{\mu\nu}\]

\textbf{Independent parameter:} \(\Lambda \approx 1.1 \times 10^{-52}\)
m\(^{-2}\)

\paragraph{Derivation of the Observed Value (Geometrical
Method)}\label{derivation-of-the-observed-value-geometrical-method}

In the Einstein equation, using the energy density parameter
\(\Omega_\Lambda \approx 0.685\) and the Hubble constant
\(H_0 \approx 67.4\) km/s/Mpc obtained from observations of the
accelerating expansion of the universe (Type Ia supernovae, CMB, BAO),

\[\Lambda = 3 \Omega_\Lambda \frac{H_0^2}{c^2} \approx 3 \times 0.685 \times \frac{(2.18 \times 10^{-18} \; \text{s}^{-1})^2}{(3.0 \times 10^8 \; \text{m/s})^2} \approx 1.1 \times 10^{-52} \; \text{m}^{-2}\]

However, this derivation is \textbf{circular}. \(H_0\) and
\(\Omega_\Lambda\) are values obtained from the observed curvature of
the universe, and using them to determine \(\Lambda\) is merely
substitution. We are not ``deriving'' the value of \(\Lambda\) from
geometry but \textbf{merely assigning the name \(\Lambda\) to the
observed curvature}.

\paragraph{Quantum Field Theory Theoretical
Prediction}\label{quantum-field-theory-theoretical-prediction}

In quantum field theory, the vacuum possesses a finite energy density as
the sum of zero-point energies of all field modes. Integrating the
zero-point energy \(\frac{1}{2}\hbar\omega\) of each mode up to an
ultraviolet cutoff \(k_{\max}\), the vacuum energy density is

\[\rho_{\text{vac}} = \frac{1}{(2\pi)^3} \int_0^{k_{\max}} \frac{1}{2} \hbar \omega_k \cdot 4\pi k^2 \, dk = \frac{\hbar}{4\pi^2 c} \int_0^{k_{\max}} \frac{1}{2} \sqrt{k^2 c^2 + m^2 c^4 / \hbar^2} \cdot k^2 \, dk\]

For a massless field (\(m = 0\)) with cutoff at the Planck scale
\(k_{\max} = k_P = \sqrt{\frac{c^3}{\hbar G}}\),

\[\rho_{\text{vac}} \sim \frac{\hbar c}{16\pi^2} k_P^4 = \frac{c^7}{16\pi^2 \hbar G^2}\]

Converting to a cosmological constant
(\(\Lambda_{\text{QFT}} = \frac{8\pi G}{c^4} \rho_{\text{vac}}\)),

\[\Lambda_{\text{QFT}} \sim \frac{c^3}{2\pi \hbar G} = \frac{1}{2\pi \ell_P^2}\]

where \(\ell_P = \sqrt{\hbar G / c^3} \approx 1.616 \times 10^{-35}\) m
is the Planck length. Therefore,

\[\Lambda_{\text{QFT}} \sim \frac{1}{2\pi \times (1.616 \times 10^{-35})^2} \approx 6.1 \times 10^{68} \; \text{m}^{-2}\]

\paragraph{Discrepancy}\label{discrepancy}

{\def\LTcaptype{none} % do not increment counter
\begin{longtable}[]{@{}
  >{\raggedright\arraybackslash}p{(\linewidth - 4\tabcolsep) * \real{0.1304}}
  >{\raggedright\arraybackslash}p{(\linewidth - 4\tabcolsep) * \real{0.4783}}
  >{\raggedright\arraybackslash}p{(\linewidth - 4\tabcolsep) * \real{0.3913}}@{}}
\toprule\noalign{}
\begin{minipage}[b]{\linewidth}\raggedright
\end{minipage} & \begin{minipage}[b]{\linewidth}\raggedright
Value (m\(^{-2}\))
\end{minipage} & \begin{minipage}[b]{\linewidth}\raggedright
Derivation method
\end{minipage} \\
\midrule\noalign{}
\endhead
\bottomrule\noalign{}
\endlastfoot
Observed value & \(1.1 \times 10^{-52}\) & Observed data of accelerating
cosmic expansion (circular) \\
QFT theoretical value & \(\sim 6 \times 10^{68}\) & Planck-scale
integration of vacuum zero-point energy \\
\textbf{Discrepancy} & \textbf{\$\approx\$120 orders of magnitude} & \\
\end{longtable}
}

\textbf{Structural problems:}

\begin{enumerate}
\def\labelenumi{\arabic{enumi}.}
\tightlist
\item
  \textbf{The 120 orders of magnitude problem.} As shown above, the
  theoretical prediction from quantum field theory vacuum energy
  deviates from the observed value by approximately 120 orders of
  magnitude. The worst theoretical prediction in all of physics. The
  essence of this discrepancy lies in quantum field theory's assumption
  that ``all modes possess zero-point energy up to the Planck scale.''
  Whether this assumption itself is correct has not been verified.
\item
  \textbf{Dimension is curvature.} \([\Lambda] = [L^{-2}]\),
  representing intrinsic curvature of spacetime. Same geometrical
  category as G (\(M\) \(\leftrightarrow\) curvature conversion).
\item
  \textbf{Why not zero?} Many physicists considered \(\Lambda = 0\) to
  be ``natural,'' but the accelerating expansion of the universe was
  discovered in 1998, confirming \(\Lambda > 0\).
\item
  \textbf{Cutoff dependence.} The QFT theoretical value is proportional
  to the 4th power of the cutoff \(k_{\max}\). There is no physical
  basis for choosing the Planck scale, and changing \(k_{\max}\) changes
  the theoretical value arbitrarily. That is, the 120-order discrepancy
  does not mean ``the theoretical prediction was wrong'' but indicates
  that \textbf{the formulation itself is incomplete}.
\end{enumerate}

\subsubsection{\texorpdfstring{5.3 QCD Vacuum Angle
\(\theta_{QCD}\)}{5.3 QCD Vacuum Angle \textbackslash theta\_\{QCD\}}}\label{qcd-vacuum-angle-theta_qcd}

\textbf{Independent parameter:} \(\theta_{QCD} < 10^{-10}\)

Experimentally nearly zero, but there is no theoretical necessity for it
to be zero. Why it is near zero is unsolved as the ``strong CP
problem.'' If a Peccei-Quinn symmetry exists, it would be structurally
zero, in which case the axion particle is predicted (not yet
discovered).

\begin{center}\rule{0.5\linewidth}{0.5pt}\end{center}

\subsection{6. Classification Results}\label{classification-results}

\subsubsection{6.1 Comparison of Formulation Levels of the Four
Forces}\label{comparison-of-formulation-levels-of-the-four-forces}

Our analysis reveals that the four forces differ greatly in their
theoretical completeness.

{\def\LTcaptype{none} % do not increment counter
\begin{longtable}[]{@{}
  >{\raggedright\arraybackslash}p{(\linewidth - 8\tabcolsep) * \real{0.1053}}
  >{\raggedright\arraybackslash}p{(\linewidth - 8\tabcolsep) * \real{0.1579}}
  >{\raggedright\arraybackslash}p{(\linewidth - 8\tabcolsep) * \real{0.2982}}
  >{\raggedright\arraybackslash}p{(\linewidth - 8\tabcolsep) * \real{0.2105}}
  >{\raggedright\arraybackslash}p{(\linewidth - 8\tabcolsep) * \real{0.2281}}@{}}
\toprule\noalign{}
\begin{minipage}[b]{\linewidth}\raggedright
Item
\end{minipage} & \begin{minipage}[b]{\linewidth}\raggedright
Gravity
\end{minipage} & \begin{minipage}[b]{\linewidth}\raggedright
Electromagnetism
\end{minipage} & \begin{minipage}[b]{\linewidth}\raggedright
Weak force
\end{minipage} & \begin{minipage}[b]{\linewidth}\raggedright
Strong force
\end{minipage} \\
\midrule\noalign{}
\endhead
\bottomrule\noalign{}
\endlastfoot
Spatial dependence & \(1/r^2\) (exact) & \(1/r^2\) (exact) &
\(e^{-Mr}/r\) (Yukawa type) & Logarithmic (1-loop approx.) \\
Computable at all distances? & \(\checkmark\) & \(\checkmark\) &
\(\checkmark\) & \(\times\) (breaks down at low energy) \\
Exact solution exists? & \(\checkmark\) & \(\checkmark\) &
\(\checkmark\) & \(\times\) (Millennium Prize Problem) \\
Dimension of charge & \(M\) (continuous) & \(Q\) (continuous, quantized)
& Discrete (\(\pm 1/2\)) & Discrete (3 colors) \\
Discretization of charge explained? & Not addressed & Unsolved (Dirac
condition) & Assumed (SU(2)) & Assumed (SU(3)) \\
Dimensional conversion constant & G (explicit) & e, \(\mu_0\epsilon_0\)
(explicit) & \(v\) (treated separately) & \(\Lambda_{QCD}\) (hidden) \\
Number of independent parameters & G, \(\Lambda\) (2) & e (1) & \(v\),
\(\theta_W\) (2) & \(\Lambda_{QCD}\) (1) \\
\end{longtable}
}

\subsubsection{6.2 Parameter Dependency
Relationships}\label{parameter-dependency-relationships}

Below we show the dependency relationships among the approximately 32
parameters. Arrows mean ``B is derived from A.''

\begin{verbatim}
§1.0 Fundamental Constants (Independent)
├── c (P1) ─────────────────┐
├── ℏ (P2) ──────────────┐  │
├── k_B (P3) ────────    │  │
├── e (P4) ──────────┐   │  │
├── G (P5) ──────┐   │   │  │
└── μ₀ε₀ (P6) ──┤   │   │  │
                 │   │   │  │
             ┌───┘   │   │  │
             ▼       ▼   ▼  ▼
          α_G ←── m_i   α = e²/(4πε₀ℏc)  ← Dependent (not independent)
                  ▲
                  │
          y_i × v/√2 = m_i
          ▲       ▲
          │       │
    Yukawa        Higgs VEV
    y_i (9)       v (1)  ← Independent parameters
                  │
                  ├── α_w = f(v, θ_W)  ← Dependent
                  ├── G_F = 1/(√2 v²)  ← Dependent
                  └── m_H = √(2λ)·v    ← Dependent
                       ▲
                       │
                  λ (1)  ← Independent parameter

Other Independent Parameters:
├── Λ_QCD (1) ── → α_s(Q²) = f(Λ_QCD, Q)  ← Dependent
├── θ_W (1)
├── Λ (1) ── Cosmological constant
├── θ_QCD (1)
├── CKM: θ₁₂, θ₂₃, θ₁₃, δ (4)
└── PMNS: Δm²₂₁, Δm²₃₂, θ₁₂, θ₂₃, θ₁₃, δ, (α₁, α₂) (7+)
\end{verbatim}

\subsubsection{6.3 TOE KPI (Acceptance Test
Specification)}\label{toe-kpi-acceptance-test-specification}

We define the evaluation criteria for a Theory of Everything as follows.

\textbf{Test cases:} Approximately 32 independent parameters

{\def\LTcaptype{none} % do not increment counter
\begin{longtable}[]{@{}
  >{\raggedright\arraybackslash}p{(\linewidth - 6\tabcolsep) * \real{0.1250}}
  >{\raggedright\arraybackslash}p{(\linewidth - 6\tabcolsep) * \real{0.3125}}
  >{\raggedright\arraybackslash}p{(\linewidth - 6\tabcolsep) * \real{0.3438}}
  >{\raggedright\arraybackslash}p{(\linewidth - 6\tabcolsep) * \real{0.2188}}@{}}
\toprule\noalign{}
\begin{minipage}[b]{\linewidth}\raggedright
ID
\end{minipage} & \begin{minipage}[b]{\linewidth}\raggedright
Category
\end{minipage} & \begin{minipage}[b]{\linewidth}\raggedright
Parameters
\end{minipage} & \begin{minipage}[b]{\linewidth}\raggedright
Count
\end{minipage} \\
\midrule\noalign{}
\endhead
\bottomrule\noalign{}
\endlastfoot
T01--T05 & Fundamental constants & c, \(\hbar\), \(k_B\), e, G & 5 \\
T06 & Cosmological constant & \(\Lambda\) & 1 \\
T07 & QCD scale & \(\Lambda_{QCD}\) & 1 \\
T08 & Electroweak mixing & \(\theta_W\) & 1 \\
T09--T17 & Yukawa couplings &
\(y_e, y_\mu, y_\tau, y_u, y_d, y_c, y_s, y_t, y_b\) & 9 \\
T18--T19 & Higgs & \(\lambda\), \(v\) & 2 \\
T20--T23 & CKM matrix &
\(\theta_{12}, \theta_{23}, \theta_{13}, \delta_{CKM}\) & 4 \\
T24 & QCD vacuum & \(\theta_{QCD}\) & 1 \\
T25--T31 & Neutrino &
\(\Delta m^2_{21}, \Delta m^2_{32}, \theta_{12}, \theta_{23}, \theta_{13}, \delta_{PMNS}, (\alpha_1, \alpha_2)\)
& 7+ \\
& \textbf{Total} & & \textbf{\$\approx\$32} \\
\end{longtable}
}

\textbf{Evaluation method:}

\[\text{KPI} = \frac{\text{Number of values derived by the theory matching experiment within significant digits}}{32}\]

\textbf{Evaluation scale:}

{\def\LTcaptype{none} % do not increment counter
\begin{longtable}[]{@{}
  >{\raggedright\arraybackslash}p{(\linewidth - 2\tabcolsep) * \real{0.3889}}
  >{\raggedright\arraybackslash}p{(\linewidth - 2\tabcolsep) * \real{0.6111}}@{}}
\toprule\noalign{}
\begin{minipage}[b]{\linewidth}\raggedright
Score
\end{minipage} & \begin{minipage}[b]{\linewidth}\raggedright
Assessment
\end{minipage} \\
\midrule\noalign{}
\endhead
\bottomrule\noalign{}
\endlastfoot
0/32 & Current Standard Model (merely returns inputs) \\
1--5/32 & Partial success (e.g., explaining \(\theta_{QCD} = 0\),
unification of coupling constants) \\
6--15/32 & Significant progress (e.g., explaining Yukawa coupling
hierarchy) \\
16--25/32 & Grand Unified Theory level \\
26--31/32 & Theory of Everything candidate \\
\textbf{32/32} & \textbf{Completion of the Theory of Everything} \\
\end{longtable}
}

\textbf{Bonus by number of input parameters:}

{\def\LTcaptype{none} % do not increment counter
\begin{longtable}[]{@{}ll@{}}
\toprule\noalign{}
Inputs & Assessment \\
\midrule\noalign{}
\endhead
\bottomrule\noalign{}
\endlastfoot
32 (all inputs) & Nothing derived (current status) \\
10--31 & Partial reduction \\
3--9 & Significant reduction \\
1--2 & Excellent TOE candidate \\
\textbf{0} & \textbf{Ultimate TOE (no initial conditions needed)} \\
\end{longtable}
}

\begin{center}\rule{0.5\linewidth}{0.5pt}\end{center}

\subsection{7. Conclusions}\label{conclusions}

\subsubsection{7.1 Reassessment of the ``19 Arbitrary
Parameters''}\label{reassessment-of-the-19-arbitrary-parameters}

Our analysis reveals the following:

\begin{enumerate}
\def\labelenumi{\arabic{enumi}.}
\item
  \textbf{The conventional ``19'' is inaccurate.} Both double-counting
  (\(\alpha\), \(\alpha_G\), \(\alpha_w\), particle masses, \(m_H\)) and
  under-counting (G, \(\Lambda\), \(\Lambda_{QCD}\), \(\theta_W\),
  Yukawa couplings, \(\lambda\), neutrinos) exist. The accurate count is
  approximately 32 independent parameters.
\item
  \textbf{Exclusion of fundamental physical constants is unjustified.}
  c, \(\hbar\), \(k_B\), e, G are excluded as ``constants,'' but the
  physical origin of their values is unresolved and they should be
  included as TOE verification targets.
\item
  \textbf{The four forces are not formulated at equal levels.} Gravity
  and electromagnetism possess exact analytical solutions (\(1/r^2\))
  and explicit dimensional conversion constants, while the strong force
  has only approximate solutions (logarithmic) that break down at low
  energies, and the weak force hides its dimensional conversion
  constants within the Higgs field.
\item
  \textbf{There is no mystery in \(\alpha \approx 1/137\).} It is a
  dependent quantity obtained automatically by substituting values into
  \(e^2/(4\pi\epsilon_0\hbar c)\), and the assessment that ``it would be
  amazing to derive it'' reflects an incorrectly posed question. What
  should be asked is the origin of the elementary charge \(e\).
\item
  \textbf{The ``hierarchy problem'' may be a pseudo-problem.} The four
  coupling constants differ in definition, spatial dependence, and
  charge dimensions. Comparing them at a single scale and lamenting that
  ``gravity alone is weak'' is the result of forcing quantities of
  different categories onto a common standard.
\end{enumerate}

\subsubsection{7.2 TOE KPI --- Priority by
Verifiability}\label{toe-kpi-priority-by-verifiability}

Verifying all approximately 32 independent parameters at once is
impossible at the current stage. The reason is that the theoretical
maturity of the four forces is fundamentally different (see §6.1).

\paragraph{Verifiable Domain: Unification of Gravity and
Electromagnetism}\label{verifiable-domain-unification-of-gravity-and-electromagnetism}

Gravity and electromagnetism satisfy the conditions necessary for
verification in the following respects:

\begin{itemize}
\tightlist
\item
  Both possess exact analytical solutions (\(1/r^2\))
\item
  The dimensions of charge are explicitly defined (mass \(M\), charge
  \(Q\))
\item
  Dimensional conversion constants are explicit (\(G\), \(e\),
  \(\mu_0\epsilon_0\))
\item
  The independent parameters involved are identified (\(c\), \(\hbar\),
  \(G\), \(e\), \(\mu_0\epsilon_0\), \(\Lambda\))
\end{itemize}

Therefore, the parameter relationships that a unified theory should
derive are clear, and pass/fail judgment through comparison with
experimental values is possible.

\textbf{KPIs for Gravity + Electromagnetism (Verifiable):}

{\def\LTcaptype{none} % do not increment counter
\begin{longtable}[]{@{}
  >{\raggedright\arraybackslash}p{(\linewidth - 4\tabcolsep) * \real{0.1081}}
  >{\raggedright\arraybackslash}p{(\linewidth - 4\tabcolsep) * \real{0.4865}}
  >{\raggedright\arraybackslash}p{(\linewidth - 4\tabcolsep) * \real{0.4054}}@{}}
\toprule\noalign{}
\begin{minipage}[b]{\linewidth}\raggedright
ID
\end{minipage} & \begin{minipage}[b]{\linewidth}\raggedright
Verification item
\end{minipage} & \begin{minipage}[b]{\linewidth}\raggedright
Pass criterion
\end{minipage} \\
\midrule\noalign{}
\endhead
\bottomrule\noalign{}
\endlastfoot
T01 & Relationship between \(G\) and \(e\) & Derive one from the other
and match experimental value \\
T02 & Value of \(\alpha = e^2/(4\pi\epsilon_0\hbar c)\) & Derive \(e\)
from structure and match \(\alpha \approx 1/137.036\) \\
T03 & Value of \(\Lambda\) & Derive
\(\Lambda \approx 1.1 \times 10^{-52}\) m\(^{-2}\) from theory \\
T04 & Relationship between mass \(M\) and charge \(Q\) & Derive the
charge quantization condition \\
\end{longtable}
}

\paragraph{Domain Difficult to Verify at Present: Weak and Strong
Forces}\label{domain-difficult-to-verify-at-present-weak-and-strong-forces}

For the weak and strong forces, the verification criterion for ``having
been unified'' is itself unclear.

\begin{itemize}
\tightlist
\item
  \textbf{Strong force:} No exact solution exists (Millennium Prize
  Problem). Perturbation theory breaks down at low energy, and the
  spatial dependence of the coupling constant \(\alpha_s\) is described
  only in 1-loop approximation. The ``correct answer'' against which to
  compare the output of a unified theory exists only as an
  approximation.
\item
  \textbf{Weak force:} The charge (weak isospin) is a discrete label
  with no SI dimension. The dimensional conversion constant is hidden in
  the Higgs field vacuum expectation value \(v\), and the separation of
  independent parameters is itself ambiguous.
\end{itemize}

Claiming to have ``unified'' incomplete theories merely means
\textbf{bundling unformulated things while they remain unformulated},
and there is no way to verify it.

\textbf{Conclusion:} The first verification step for a Theory of
Everything is the construction of a unified theory of gravity and
electromagnetism, and the verification of T01--T04 above. Unification of
the weak and strong forces becomes a verification target only after the
rigorous formulation of each theory individually is completed.

\subsubsection{7.3 Limitations of This
Paper}\label{limitations-of-this-paper}

This paper has classified existing parameters and organized the
problems, but has not derived the parameters themselves. Structural
derivation of each parameter becomes possible only at the stage of
implementing the structural requirements formulated in the prior work
{[}1{]} into concrete geometry.

\begin{center}\rule{0.5\linewidth}{0.5pt}\end{center}

\subsection{References}\label{references}

{[}1{]} N. Kihara, ``Formulation of Structural Requirements That a
Theory of Everything Must Satisfy,'' 2026. DOI:
\href{https://doi.org/10.5281/zenodo.19601592}{10.5281/zenodo.19601592}

{[}2{]} Particle Data Group, ``Review of Particle Physics,'' \emph{Phys.
Rev.~D} \textbf{110}, 030001 (2024).

{[}3{]} S. Weinberg, ``A Model of Leptons,'' \emph{Phys. Rev.~Lett.}
\textbf{19}, 1264 (1967).

{[}4{]} P. W. Higgs, ``Broken Symmetries and the Masses of Gauge
Bosons,'' \emph{Phys. Rev.~Lett.} \textbf{13}, 508 (1964).

{[}5{]} G. 't Hooft, ``Naturalness, Chiral Symmetry, and Spontaneous
Chiral Symmetry Breaking,'' \emph{NATO Adv. Study Inst. Ser. B Phys.}
\textbf{59}, 135 (1980).

{[}6{]} A. Einstein, ``Die Feldgleichungen der Gravitation,''
\emph{Sitzungsberichte der Preussischen Akademie der Wissenschaften zu
Berlin}, 844--847 (1915).

{[}7{]} D. J. Gross and F. Wilczek, ``Ultraviolet Behavior of
Non-Abelian Gauge Theories,'' \emph{Phys. Rev.~Lett.} \textbf{30}, 1343
(1973).

{[}8{]} Planck Collaboration, ``Planck 2018 results. VI. Cosmological
parameters,'' \emph{Astron. Astrophys.} \textbf{641}, A6 (2020).

\end{document}
